\documentclass[11pt]{article}
\usepackage[margin=.85in]{geometry}
\usepackage{amsmath,amssymb,booktabs,hyperref}
\hypersetup{colorlinks=true,urlcolor=blue}
\newcommand{\Q}{\mathbb Q}
\newcommand{\Z}{\mathbb Z}
\title{\vspace{-1.7em}Seven terms are necessary at discriminant \(-152\)}
\author{A supplement to the ProveIt GRH article}
\date{27 November 2026}
\begin{document}
\maketitle
\begin{abstract}
The merged ProveIt article explicitly asks whether the least support
of a rational polynomial annihilating a singular modulus of discriminant
\(-152\) is seven. The answer is yes: every six distinct integer powers
of the original singular modulus are linearly independent.
The same method determines the exact least support for thirteen further
orders. These deductions are unconditional and are not Lean-formalized.
\end{abstract}

\section{The certificate for \(-152\)}
Let \(H_\Delta\) be the monic ring class polynomial of discriminant
\(\Delta\). Its roots are the singular moduli of that order, and it is
the minimal polynomial of each. Write \(s(\Delta)\) for the least
number of nonzero terms in a nonzero rational polynomial vanishing
at such a root.

The \href{https://github.com/VladimirReshetnikov/ProveIt/blob/e7f7a226563fb14c6f0b796322e43b048c548953/Analysis/GeneralizedRiemannHypothesis/Article/sections/smd.tex#L1552-L1563}{article's six-terms question}
asks in particular for \(s(-152)\) and \(s(-87)\). The prior
theorem on powers covers the \(182\)nd power of the \(-152\) modulus,
not a relation with arbitrary exponents in the original modulus.
For the original modulus, its monic degree-six polynomial satisfies
\[
 \frac{H_{-152}(64+256Y)}{256^6}\in\Z[Y].
\]
For transparency, the entire transformed polynomial is
\[
\begin{aligned}
Y^6&-258774476247186Y^5+259766390712739705665Y^4\\
 &+408822444341150448859180Y^3+648287850356451166578247815Y^2\\
 &-931493942482171555920209906Y+1404883983219431921036769431.
\end{aligned}
\]
Every root \(Y\) is \(2\)-adically integral: otherwise the leading
monomial would have strictly lower valuation than all other terms.
Thus every conjugate \(\alpha\) of the original singular modulus
satisfies \(v_2(\alpha/64-1)\ge2\), where \(v_p(p)=1\).

\paragraph{Why the disk excludes a six-term relation.}
Here is the \href{https://github.com/VladimirReshetnikov/ProveIt/blob/e7f7a226563fb14c6f0b796322e43b048c548953/Analysis/GeneralizedRiemannHypothesis/Article/sections/smd.tex#L64-L109}{article's local argument}.
In a disk \(v_p(x/A-1)\ge m>1/(p-1)\), substitute
\(x=A\exp(p^mY)\) for the integer \(m\)'s used here; the inverse uses
the convergent logarithm and the disc becomes \(v_p(Y)\ge0\).
If a \(t\)-term polynomial \(P=\sum_{j=1}^{t}c_jX^{n_j}\) is evaluated
there, its \(k\)-th power-series coefficient is \(M_kp^{mk}/k!\) with
\(M_k=\sum_j c_jA^{n_j}n_j^k\). The first \(t\) moments are not all
zero by the Vandermonde determinant. All later moments satisfy the
monic recurrence with polynomial \(\prod_j(T-n_j)\in\Z[T]\), so
their valuations are bounded below by the minimum of the first \(t\).
\newpage
Consequently, if \(w(k)=mk-v_p(k!)\) satisfies
\[
 \inf_{k\ge h}w(k)>\max_{i<h}w(i),
\]
Strassmann's theorem gives at most \(h-1\) zeros when \(t\le h\):
all coefficients from index \(h\) onward have larger valuation than
some earlier coefficient. Convergence follows from Legendre's
factorial bound. Here \(p=2,m=2,h=6\). The maximum for \(i<6\)
is \(7\), whereas \(w(k)=k+S_2(k)\ge8\) for \(k\ge6\), with
\(k=6\) checked directly. Any rational relation at one root would
vanish at all six by minimality, contradicting the bound. The
class polynomial provides a relation with at most seven terms;
therefore \(s(-152)=7\). Multiplying by a common power allows
negative exponents as well. The result is about arbitrary exponents,
not just exponents in one residue class.

\section{Thirteen more exact values}
For each row set \(A=p^v u,\ B=p^{v+m}\). The exact check is that
\(H_{-D}(A+BY)/B^h\) is monic integral, which puts all \(h\) roots in
the depth-\(m\) disk about \(A\). The last column gives the valuations
of the coefficients of \(p^{-vh}H_{-D}(p^v(u+Z))\), in ascending order.
They are at least \(m(h-i)\) at index \(i\), as required.
\begin{center}
\small
\begin{tabular}{rrrrrl}
\toprule
\(D\) & \(h\) & \(p\) & \(v\) & \((u,m)\) & Coefficient valuations\\
\midrule
152  &6&2&6&(1,2)&12, 11, 8, 8, 4, 3, 0\\
244  &6&2&6&(3,2)&12, 13, 8, 10, 4, 5, 0\\
300  &6&2&4&(3,2)&12, 11, 8, 8, 4, 3, 0\\
451  &6&3&3&(1,2)&12, 12, 9, 6, 8, 3, 0\\
472  &6&2&6&(1,2)&12, 11, 8, 8, 4, 3, 0\\
588  &6&2&4&(3,2)&12, 11, 8, 9, 4, 3, 0\\
628  &6&2&6&(3,2)&12, 13, 8, 10, 4, 5, 0\\
859  &7&3&3&(1,1)&14, 12, 11, 8, 6, 5, 2, 0\\
867  &6&3&3&(2,2)&14, 12, 9, 7, 4, 5, 0\\
1083 &6&17&0&(8,1)&6, 5, 4, 4, 2, 1, 0\\
1147 &6&3&3&(1,2)&12, 11, 9, 6, 5, 3, 0\\
1483 &7&3&3&(1,1)&14, 12, 11, 8, 6, 5, 2, 0\\
2467 &7&3&3&(1,1)&14, 12, 11, 8, 6, 5, 2, 0\\
2707 &7&3&3&(1,1)&14, 12, 11, 8, 6, 5, 2, 0\\
\bottomrule
\end{tabular}
\end{center}
For \((p,m,h)=(3,2,6)\) the initial maximum is \(9\), and for
\((17,1,6)\) it is \(5\). Legendre's lower bound
\(w(k)\ge mk-(k-1)/(p-1)\) exceeds the respective maxima at \(k=6\)
and increases afterward. For \((3,1,7)\) the maximum is \(4\),
\(w(7)=5\), and the same lower bound works from \(k=8\). Thus all
rows give \(s(-D)=h+1\). This does not settle \(-87\) or the universal
six-terms question.
\newpage
\section{Independent replay of the inputs}
The included vectors were taken from the
\href{https://github.com/VladimirReshetnikov/ProveIt/blob/10802a6b36/Analysis/GeneralizedRiemannHypothesis/Research/grh_sparse_singular_moduli/checks/five_small_witnesses.json}{pinned coefficient fixture}.
The exact checker uses both the binomial theorem and Horner composition
on the affine substitution. For the factorial comparison it checks
integers until the linear lower bound from Legendre becomes strict.
The two substitution methods agreed on all 102 coefficients. The
checker also enumerates all primitive reduced forms of each listed
discriminant, independently establishing the degrees.

A second script verifies the polynomial inputs directly using complex
interval arithmetic. For a reduced form \((a,b,c)\) of discriminant
\(-D\), use \(q=\exp(2\pi i\tau)\), with
\(\tau=(-b+i\sqrt D)/(2a)\). The classical modular identity is
\[
 j(\tau)=\frac{(1+240\sum_{n\ge1}\sigma_3(n)q^n)^3}
                  {q\prod_{n\ge1}(1-q^n)^{24}}.
\]
The uniform bound \(|q|\le e^{-\pi\sqrt3}<r=1/200\) comes from
reduction. At cutoff \(N=112\) the series tail in the numerator's
base is bounded by
\(240(N+1)^4r^{N+1}/(1-16r)\), since \(\sigma_3(n)\le n^4\).
If \(V=\prod_{n>N}(1-q^n)^{24}\), then the sum of logarithms has
absolute value at most \(24r^{N+1}/((1-r)(1-r^{N+1}))\).
In particular \( |V^{-1}-1|\le100r^{N+1}\).
The script encloses these errors in rectangles, multiplies the
enclosures of \(X-j(\tau)\), and verifies that each coefficient
interval contains exactly the asserted integer on the real axis.
Since the ring class polynomial has integer coefficients, this
independently recovers all 102 input integers. The replay uses
\texttt{mpmath.iv} interval arithmetic, not PARI/GP. This is an
input audit plus a written proof, not a formal verification of
complex interval arithmetic or a statement about publication priority.
\end{document}
